\UseRawInputEncoding
\documentclass[10pt, a4paper]{article}
%\documentclass[12pt, a4paper, twoside]{book}

\usepackage{helvet}
\usepackage{hyperref}
\usepackage{graphics}
\usepackage{graphicx}
\usepackage{listings}
\usepackage{textcomp}
\usepackage[
	a4paper,
	outer=2cm,
	inner=4cm,
	top=2cm,
	bottom=2cm
]{geometry}
\usepackage{float}
\usepackage{tabularx}
\usepackage{xcolor}
%\usepackage[disable]{todonotes}
%\usepackage{color, soul}
\usepackage{amsmath}
%\usepackage{algorithmicx}
%\usepackage[noend]{algpseudocode}
%\usepackage{algorithm}
\usepackage{framed}
\usepackage{subcaption}
%\usepackage{titlepic}
%\usepackage{fancyhdr}
%\usepackage[simplified]{styles/pgf-umlcd}
\usepackage{shorttoc}
\usepackage{url}
%\usepackage{paralist}
\usepackage{dirtytalk}
%\usepackage{verbatim}

\definecolor{grey}{rgb}{0.9, 0.9, 0.9}
\definecolor{dkgreen}{rgb}{0,0.6,0}
\definecolor{dkred}{rgb}{0.6,0,0.0}
\lstdefinestyle{DOS}
{
    backgroundcolor=\color{black},
    basicstyle=\scriptsize\color{white}\ttfamily,
    stringstyle=\color{white},
    keywords={}
}

\lstdefinestyle{makefile}
{
    numberblanklines=false,
    language=make,
    tabsize=4,
    keywordstyle=\color{red},
    identifierstyle= %plain identifiers for make
}

\lstset{
  language=Java,                % the language of the code
  basicstyle=\footnotesize\ttfamily,
  numbers=left,                   % where to put the line-numbers
  stepnumber=1,                   % the step between two line-numbers. If it's 1, each line
  numbersep=5pt,                  % how far the line-numbers are from the code
  backgroundcolor=\color{white},      % choose the background color. You must add \usepackage{color}
  showspaces=false,               % show spaces adding particular underscores
  showstringspaces=false,         % underline spaces within strings
  showtabs=false,                 % show tabs within strings adding particular underscores
  frame=single,                   % adds a frame around the code
  rulecolor=\color{black},        % if not set, the frame-color may be changed on line-breaks within not-black text (e.g. comments (green here))
  tabsize=2,                      % sets default tabsize to 2 spaces
  captionpos=b,                   % sets the caption-position to bottom
  breaklines=true,                % sets automatic line breaking
  breakatwhitespace=false,        % sets if automatic breaks should only happen at whitespace
  keywordstyle=\color{blue},          % keyword style
  commentstyle=\color{dkgreen},       % comment style
  stringstyle=\color{dkred},         % string literal style
  columns=fixed,
  extendedchars=true,
  frame=single,
}

%\renewcommand{\chaptername}{Topic}

% New definitions
%\algnewcommand\algorithmicswitch{\textbf{switch}}
%\algnewcommand\algorithmiccase{\textbf{case}}
%\algnewcommand\algorithmicassert{\texttt{assert}}
%\algnewcommand\Assert[1]{\State \algorithmicassert(#1)}%
% New "environments"
%\algdef{SE}[SWITCH]{Switch}{EndSwitch}[1]{\algorithmicswitch\ #1\ \algorithmicdo}{\algorithmicend\ \algorithmicswitch}%
%\algdef{SE}[CASE]{Case}{EndCase}[1]{\algorithmiccase\ #1}{\algorithmicend\ \algorithmiccase}%
%\algtext*{EndSwitch}%
%\algtext*{EndCase}%

%\pagestyle{fancy}
%\fancyhf{}
%\fancyhead[RO, LE]{\small \rightmark}
%\fancyfoot[RO, LE]{\small \thepage}

\begin{document}

%\frontmatter

\begin{titlepage}
\vspace*{5cm}
\begin{center}
\includegraphics[width=.5\textwidth]{images/EdNapUniLogoCMYK.pdf}~\\[1cm]

\textsc{\Large Edinburgh Napier University}\\[1.5cm]

\textsc{\LARGE \bfseries SET08101 Web Tech}\\[0.5cm]

\hrulefill \\[0.4cm]
{\huge \bfseries Lab 6 - Design Hacks for Coders \\[0.4cm] }
\hrulefill \\[1.5cm]

\begin{minipage}{0.4\textwidth}
\begin{flushleft} \large
\textbf{Dr Simon Wells} \\
\end{flushleft}
\end{minipage}

\vfill

\end{center}
\end{titlepage}

%\shorttoc{Overview}{0}

%\setcounter{tocdepth}{2}
%\cleardoublepage
%\tableofcontents
%\listoffigures
%\listofalgorithms
%\addtocontents{toc}{~\hfill\textbf{Page}\par}

%\mainmatter

%\input{sections/labs/04_ui}

\section{Aims}
\paragraph{} By the end of the practical portion of this topic you will:

\begin{itemize}
\item Have developed some idea of the range of tools and techniques available for prototyping both web pages and the network of relationships between pages in a site that you're designing.
\item Have used placeholder text generation tools to help you lay out portions of pages before the final content is ready.
\item Have used colour tools to select a palette of visually appealing and accessible colours for use on a site.
\item Have investigated and applied style guides to the design of a website.
\end{itemize}

\begin{framed}
{\bf{NOTICE:}  Some of the exercises today build upon learning from previous weeks. Importantly, you should be noticing that many of the techniques and concepts we have studied so far will draw together to help you tackle the first coursework assignment.}
\end{framed}


\paragraph{} Many coders say that they ``can't do design'' but that is untrue. You don't need to be a graphic designer in order to produce solid, pleasing, and usable designs. All you really have to do is practise a few tools that give short-cuts to some of the more difficult tasks like selecting colours and to follow some simple design principles, like keeping things simple and uncluttered, rather than crowded. Do this, practise, and maintain a critical eye for your own and other's designs and you will soon be able to produce acceptable website designs. Remember, your site does not have to be the most beautiful, artistic, or novel in the world, it just has to communicate your data to your user. This is it's primary job. So you should be considering order of data, arrangement of data, linking between and within data, how data can be broken into various subsections or pages. This will all inform the design of your overall navigation and how your individual pages should be laid out. 

\section{Activities}

\subsection{Design: Prototyping with Wireframe \& Mockup Tools}
\paragraph{} There are two key activities in the early design phase for a website. Firstly, determining what pages are required for your site and how to navigate consistently between them, and secondly, determing how the content on your pages should be arranged and laid out. Neither task can be successfully completed in isolation from the other as navigation between pages can effect navigation elements of individual pages, and the amount of content shown on a given page can affect the number or range of pages required overall for the site.

\paragraph{} There are many tools that can be used to mockup a website during the design phase. The simplest is pen \& paper, used to draw out sketches of how a page should look; where the major elements of the page should be placed in relation to each ther. Similarly, a simple diagram showing how a site can be navigated is a useful design artefact that can be produced on paper. A whiteboard and a copious supply of post-it notes is also a useful \emph{analogue} toolset as you can easily make changes to your design and move things around.

\paragraph{} It is a good idea to start sketching out a navigation tree early in the design phase, even if it just means that you have some arrows, representing hyperlinks, between boxes, representing pages. Give thoses boxes a name. You know that there will be one called index.html and perhaps others such as an about page, a help page, a contact page. You should also consider the way that your content is organised, perhaps there is a hierarchy in which your pages can be grouped and then navigated, for example, a bookshop website might have a page for each book, however, in order to navigate perhaps thousands of pages for individual books, they may in turn be organised by genre, rating, price, edition, or any number of other factors. These kinds of groupings give us a way to sort a potentially huge number of pages into a more organised and managable set of groups of pages.


\paragraph{} Almost any graphics package will support you in drawing a navigation tree. However, there are also a heap of online tools that provide methods for mocking up website designs. These are some free ones:

\begin{itemize}
\item \url{https://wireframe.cc/}
\item \url{http://mockflow.com}
\item \url{https://www.invisionapp.com/}
\item \url{https://mockingbot.com/}
\end{itemize}

\paragraph{} There are many others, a search for something along the lines of ``web design wireframe tool'' (try replacing `wireframe' with `mockup') will reveal many more . You should experiment with a few and decide if the online mockup tool is something that you want to work with. Many professional developers are happy to work with paper or whiteboard, and more generic graphcis tools so you should make a decision that works for you. One thing to now is that mockups  \& wireframes are merely prototypes. There is a range of prototypes from the lowest fidelity hand-drawn sketch right through to something built with HTML \& CSS that is just lacking content and sign-off by a decision maker. Choosing the right level of fidelity at which to prototope depends upon the time available, the complexity of the problem, and the context in which the design work is being performed, for example, if you were trying to win a contract then a medium fidelity prototype might sell your vision to a client better than a wireframe outline. You will build experience and knowledge about how best to approach the prototyping phase.

\paragraph{} Try out a few mockup tools (as well as pencil \& paper) and use them to mockup your ideas for your assignment. Try out a few alternative designs for how you would lay out your page.


\subsection{Design: Mocking up Content with Placeholder Text}
\paragraph{} Until the actual copy for your site is written, a job that is often performed by another member of your team, or even someone outside your team entirely, it can be useful to mockup the content so that you have something to place on screen. This is particularly important as you start to move from early design toward implementation. Whilst you can just copy \& paste text from elsewhere or even just place graphical blocks to indicate where content should go, there are both some advantages and disadvantage. For example, copying \& pasting existing text from elsewhere can lead to your placeholder text becoming accidentally mixed with live content but there is the advantage of existing text at least having similar rhythms and cadences to your live text. Graphical blocks will ``block out'' where text should go but this will not have very many similarities with actual text. A solution it to use ``Greeked'' text. Text sections, usually in Latin, that have similar patterns in terms of word length and paragraph length to lots of modern languages but are difficult to mistake for those same languages. Greeking has been used for centuries as a method to get placeholder text to help lay out the printed page before the final text was ready. Whilst we aren't dealing with printed pages necessarily, there are sufficient similarities for this to be a useful technique\footnote{Not everyone agrees that Greeking is either good or useful. Some people think that the design should be done with the actual text that will be displayed on screen. Whilst this is true in an ideal world, there are many reasons which such an approach might be less pragmatic. Can you think of any reasons why you might not be able to use the actual text?}.

\paragraph{} The easiest way to take advantage of Greeking is to have a section of writing by Cicero ready to use, e.g. 

\begin{framed}
Lorem ipsum dolor sit amet, consectetur adipiscing elit, sed do eiusmod tempor incididunt ut labore et dolore magna aliqua. Ut enim ad minim veniam, quis nostrud exercitation ullamco laboris nisi ut aliquip ex ea commodo consequat. Duis aute irure dolor in reprehenderit in voluptate velit esse cillum dolore eu fugiat nulla pariatur. Excepteur sint occaecat cupidatat non proident, sunt in culpa qui officia deserunt mollit anim id est laborum.
\end{framed}


\paragraph{} Just copy \& past it into your design whenever you need it. Until of course, you have actual text for the site you're developing, then it might be better to use the real thing.

\paragraph{} There are many sites that provide portions of text for use as Greeking placholder text. Some of these are quite fun:

\begin{itemize}
\item Samuel L. Ipsum: \url{http://slipsum.com/}
\item Picksum Ipsum: \url{http://www.picksumipsum.co.uk/}
\item Fillerama: \url{http://fillerama.io/}
\item TV Ipsum: \url{https://tvipsum.com/}
\end{itemize}

\paragraph{} See if you can find any other sites that provide similar functionality. It's worth having a bunch of them bookmarked ready for use during a design.

\paragraph{} Consider where in your sites, both those developed in previous weeks and you upcoming assignment site, you could use some mocked up text to help fill out some of the pages. Try this out. Notice how a little placeholder text can help flesh out the design to give an idea of what it will look like. This is an important step on the journey from mockup through to initial implementation and helps you to see how the various parts of your design work together (or not).

\subsection{Design: Colour}
\paragraph{} Once you have an idea of the content for your page, how each page should look in terms of where the content is laid out in relation to other bits, you'll probably want to consider colours. We'll refer to the set of colours that you use in your site as your site's ``palette''. It's actually fairly difficult to select colours before you have at least an initial wireframe or mockup sketch because, until then, you won't have a good sense of how many elements you will have in your site and consequently, the size of the palette that you need to assemble. That's why we don't need to start with setting up a colour scheme, rather, as in this lab, we start with an idea, flesh this out a little with some mockups, then try an initial implementation, focussing on layout rather than style. Only after this do we start to consider how things should look (although in reality we might also be sketching out other ideas and developing an aesthteic sense for how things should look in parallel\footnote{This only goes to show how there is no single route to a good design, but rather an iterative path where we refine our ideas until we have something that is good enough}).

\paragraph{} Colour is an extremely important aspect of almost any form of design. Good use of colour can lead to a pleasant experience and ease of use for your user. However poor use of colour can have the opposite effect, leading not just to a poor user experience, but actually making a site nigh on or even actually unusable. For example, poor contrast between text colour and the background colour, can make reading very difficult. As a rule we probably want our text to be easy to read but there are occasions when making text more difficult to read is a design feature. for example, some social discussion sites decrease the contrast between text and background on low rated comments. A good example of this are within the discussions at Hacker News\footnote{\url{https://news.ycombinator.com/}}. Other sites get a similar effect using other approaches than colour, for example, moderators on the comment boards at Boing Boing\footnote{\url{http://www.boingboing.net}} remove the vowels from any comment that is felt not to meet community standards, a process known as \emph{disemvowelling}.

\paragraph{} An important point with colour is not to just choose your favourite colours and use them regardless. Instead, select a dominate colour to guide you, for example, when designing a site for an existing organistation there are often logos and colours that are part of the brand and identity of the organisation. You should use these as a starting point. You should then determing how many colours you require for your mocked up design, e.g. make a list of elements that need to be depicted using different colours, taking a note of where higher contract is required, and not forgetting backgrounds of elements. This should give you an idea of the size of palette that you will need. You can then use a palette tool to select colours for you that work together. For example, the following tools are currently popular online:

\begin{itemize}
\item \url{https://coolors.co/}
\item \url{http://www.colourlovers.com/palettes}
\item \url{http://paletton.com/}
\item \url{http://www.color-hex.com/color-palettes/}
\item \url{https://color.adobe.com/create/color-wheel/}
\end{itemize}

\paragraph{} As with the greeking and mockup tools, it is worth having a collection of bookmarks for these sites in your \emph{toolbox} for when you need them. Some tools are also easier to use than others, and some tools give you more features, so it is worth trying a few out.

\paragraph{} Determine the number of features that you'll need to provde colours for for your ROT13 encoder (or another site). Use this knowledge to generate some colour schemes and try them out by editing your CSS file.


\subsection{Design: Style Guides}
\paragraph{} As you start to make decisions about how your site will look, you will begin to build up a library of HTML elements and associated styles. In it's simplest form, a style guide is merely a collection of all elements and their styles that you plan to use anywhere on your site, all presented together in one place so that they can all be seen in context. If the style guide looks terrible, with clashing colours, typography, layout, etc. then it is unlikely that any single page within your site will also look good if it applies your style. So getting a style guide in place early can help save you some work later on.

\paragraph{} Take a look at some example style guides. A good place to start is this one \url{http://oli.jp/2011/style-guide/} which has all of the styles used in the site demonstrated on a single page so that you can easily see how any given new element that you want to add elsewhere in your site will look.

\paragraph{} A slightly more advanced alternative is this style guide: \url{https://paulrobertlloyd.com/styleguide/} which groups the styles of a series of pages that cover scopes, utillities, and components. A good place to learn more is the Style Guides site\footnote{\url{http://styleguides.io/}} which has links to hundreds of examples, and can also be a good place to get some ideas for effective designs.

\paragraph{} Finally, investigate some of the Global Experience Languages (GELs) that exist. A good place to start is with the BBC GEL\footnote{\url{https://www.bbc.co.uk/gel}}, Edinburgh University\footnote{\url{https://gel.ed.ac.uk/}}, NHS England Identity Guidelines\footnote{\url{https://www.england.nhs.uk/nhsidentity/}}, UK Government Digital Service Design Principles\footnote{\url{https://www.gov.uk/guidance/government-design-principles}}, and Transport for London's Design Standards\footnote{\url{https://tfl.gov.uk/info-for/suppliers-and-contractors/design-standards?intcmp=5837}}. \paragraph{} Implementing your own style guide can be as simple as creating an additional page within your site which displays all the presentational elements that you've used (or have designed ready to be used once the content is available, for example, you might have planned ahead for styling features that aren't yet implemented). Create your own style guide for your ROT13 encoder (or for another site that you've designed elsewhere) that brings together the aspects you've looked at so far.

\paragraph{} The important takeaway idea is that style guides and experience languages attempt to depict clarity of ideas, as well as documentation of design decisions. Getting both of these right can help lead toward clean, consistent, and extensible website designs.

\paragraph{} Try developing a style guide for one of the the sites that you've already built in previous labs, or else start experimenting with a guide for your assignment.

\end{document}


%\begin{framed}
%\end{framed}


%\begin{lstlisting}
%\end{lstlisting}


%\begin{figure}[H]
%\centering
%\includegraphics[width=0.75\textwidth]{images/}
%\caption{}
%\label{fig:}
%\end{figure}
